% English translation of checked_pages/esq_119_049.tex; source: cote 119, site 49, PDF page 50, typescript page 47.
% AI-translated by GPT-6 Astra (OpenAI, ChatGPT Pro mode).
\scanpage{119/49}
\mspage{47}

\uline{of a bigon}. (The first of these operations also plays a
key role in a "dynamic" theory of systems of pseudolines
in a real projective plane.) One of the first questions
that arise here is to determine the different \uline{immersion
classes} of a circle or a disk (say) modulo these elementary
operations; another is to see which immersions
of the boundary of the disk come from an immersion of the
disk, and to what extent the former determine the latter.
Here again, it seemed to me that a systematic study
of the relevant moduli varieties (of infinite dimension in this
case, unless one manages to give a purely combinatorial
version of them) should provide the most effective ``focus'',
in a sense forcing us to ask ourselves the most pertinent
questions. Unfortunately, reflection on even the
most obvious and down-to-earth questions has remained embryonic.
As the only tangible result, I can point to a
theory of canonical ``dévissage'' of a stable cellular immersion
of the circle in a surface into ``indecomposable'' immersions,
by ``telescoping'' such immersions. Unfortunately, I have not
succeeded in seeing my insights into the question turn into
a DEA research project, nor other insights (into a complete
theoretical description, in terms of fundamental groups of
topological 1-complexes, of the immersions of a surface with boundary
that extend a given immersion of its boundary) into the beginning of a
state doctorate thesis\tnote{TN-049-01}{``State doctorate'' renders ``doctorat d'état''. The later ``third-cycle theses'' renders ``thèses de troisième cycle''; the source uses these degree designations.}...

A third theme, pursued simultaneously for three
years at various teaching levels (from the elective for
first-year students to three third-cycle theses currently
being pursued on this theme), concerns the topological-combinatorial
classification of systems of lines or pseudolines.
On the whole, my students' participation here has been less
disappointing than elsewhere, and I have sometimes had the pleasure of learning
from them interesting things that would not have occurred to me.
Our joint reflection has of necessity, however, been limited

\margin{49}
